\tzpanchor{0106}
\tzpentryheader{0106}{Phase-Coherence Matrix}

\begingroup\small
\begingroup\scriptsize
\setlength{\tabcolsep}{4pt}
\renewcommand{\arraystretch}{0.92}
\noindent\begin{tabularx}{\linewidth}{@{}p{0.24\linewidth}p{0.36\linewidth}X@{}}
\textbf{UUID} & \multicolumn{2}{l@{}}{\tzpcode{7b900d6a-c9e4-5927-81ce-653d492a8883}}\\
\textbf{PiID} & \tzpcode{8651328230664} & \textbf{TZPi}\quad \tzpcode{8}\quad \textbf{PiPE}\quad \tzpcode{113}\\
\textbf{RTEID} & \textbf{Poloidal $\theta$}\quad \tzpcode{411.6215 rad} & \textbf{Toroidal $\phi$}\quad \tzpcode{00.6782 rad}\\
\textbf{TZPID} & \tzpcode{ID0106} & \textbf{Node Dependencies}\quad \tzpidref{0105}\\
\textbf{Registry} & \tzpcode{registry\_id\_0106} & \textbf{Scale}\quad coherence lattice\\
\textbf{LOC Classification} & QA188 matrix analysis & QC174.45 phase coherence\\
\textbf{arXiv Classification} & Primary: cond-mat.stat-mech \quad Cross-list: quant-ph, math-ph & \\
\textbf{Primary Support} & Pairwise Coherence Entry & Spectral Weighting\\
\textbf{Closure Support} & Correlation Kernel & \\
\end{tabularx}\par
\endgroup
\endgroup

\tzpsection{Canonical Statement}
This matrix encodes how phases on a helical lattice are correlated.

\tzpsection{Canonical Equation}
\[
C_{ij} = \left\langle e^{i(\theta_i-\theta_j)} \right\rangle e^{-\lambda_i/\Lambda}
\]
\[
G(x)= e^{-\lvert x\rvert/\xi}\cos(k_F x)
\]

\begin{minipage}[t]{0.49\linewidth}
\tzpsection{Canonical Symbol Key}
\begingroup
\small
\raggedright
\textbf{Source equation symbols.}\par
\begin{itemize}[leftmargin=1.35em,itemsep=1pt,topsep=1pt,parsep=0pt]
    \item $C_{ij}$: source equation symbol.
    \item $\theta_i$: source equation symbol.
    \item $\theta_j$: source equation symbol.
    \item $\lambda_i$: source equation symbol.
    \item $\Lambda$: lemniscatic representative or canonical branch object.
    \item $G$: source equation symbol.
    \item $k_F$: source equation symbol.
\end{itemize}
\endgroup
\end{minipage}\hfill
\begin{minipage}[t]{0.47\linewidth}
\tzpsection{Wolfram Form}
\begingroup
\scriptsize
\raggedright
\textbf{Source Wolfram model.}\par
\begin{Verbatim}[fontsize=\tiny,breaklines=true,breakanywhere=true]
(* TZPID ID0106 generated source Wolfram model *)
ClearAll[K0106, Cr0106, T, R, A, W, I, N];
eq0106$1 = HoldForm[CIj == Inner[Exp[I*(thetaI - thetaJ)]]*Exp[-lambdaI/Lambda]];
eq0106$2 = HoldForm[G[x] == Exp[-Abs[x]/xi]*Cos[kF*x]];
eq0106$3 = HoldForm[CIj == Inner[Exp[I*(thetaI - thetaJ)]]];

K0106 = HoldForm[{eq0106$1, eq0106$2, eq0106$3}];

N = "normalized TRAWIN closure object for Phase-Coherence Matrix";
I = "Correlation Kernel integration/comparison layer";
W = "Spectral Weighting wave or operator carrier";
A = "Pairwise Coherence Entry admissible action/amplitude condition";
R = "Spectral Weighting rotation/curl preservation layer";
T = "Pairwise Coherence Entry local source condition";

Cr0106[expr_] := HoldForm[N[I[W[A[R[T[expr]]]]]]];

Result0106 = Cr0106[K0106];
\end{Verbatim}
\endgroup
\end{minipage}

\par\vspace{0.45\baselineskip}
\noindent\begin{minipage}[t]{\linewidth}
\begingroup
\small
\raggedright
\textbf{TRAWIN Operator Key.}\par
\noindent\textbf{Time/localization operator:} $\mathsf{T}\!\left[\mathrm{local}\right]$ Pairwise Coherence Entry local source condition.\par
\noindent\textbf{Rotation/curl operator:} $\mathsf{R}\!\left[\nabla\times\right]$ Spectral Weighting rotation/curl preservation layer.\par
\noindent\textbf{Action/amplitude operator:} $\mathsf{A}\!\left[\mathrm{admissible}\right]$ Pairwise Coherence Entry admissible action/amplitude condition.\par
\noindent\textbf{Wave/d'Alembertian operator:} $\mathsf{W}\!\left[\Box\right]$ Spectral Weighting wave or operator carrier.\par
\noindent\textbf{Integration operator:} $\mathsf{I}\!\left[\int\right]$ Correlation Kernel integration/comparison layer.\par
\noindent\textbf{Normalization operator:} $\mathsf{N}\!\left[\mathrm{normalized}\right]$ normalized TRAWIN closure object for Phase-Coherence Matrix.\par
\endgroup
\end{minipage}

\clearpage
\noindent\textbf{[ID 0106] Phase-Coherence Matrix}\par
\vspace{0.35em}
\tzpsection{Derivation Layer}
\paragraph{Primary Equation.}
\[
C_{ij} = \left\langle e^{i(\theta_i-\theta_j)} \right\rangle e^{-\lambda_i/\Lambda}
\]
\[
G(x)= e^{-\lvert x\rvert/\xi}\cos(k_F x)
\]

\paragraph{Primary Derivation.}
The relation is obtained by differentiating and applying the relevant definitions. yielding the required identity.
\paragraph{Equation Type.}
This statement is classified as a other relation.

\paragraph{Interpretation.}
The identity governs the objects appearing in the equation. In the TRAWIN reading, the equation functions as the generalized carrier for the entry’s information content. Within the CHL view, the derivation provides a witness morphism connecting the canonical proposition to its proof certificate.

\par\vspace{0.35em}
\noindent\textbf{Derivable Consequences.}\quad
\begingroup
\footnotesize
\raggedright
The canonical equation anchors Phase-Coherence Matrix by connecting the source condition to the derivation layer and proof certificate.
\par\endgroup

\tzpsection{Equation Support}
\noindent\textbf{1. Pairwise Coherence Entry}\par
\[
C_{ij} = \left\langle e^{i(\theta_i-\theta_j)} \right\rangle
\]
\noindent This introduces synchronization as a measurable comparison between phase sectors rather than a single global scalar.\par
\noindent\textbf{2. Spectral Weighting}\par
\[
C_{ij} = \left\langle e^{i(\theta_i-\theta_j)} \right\rangle e^{-\lambda_i/\Lambda}
\]
\noindent This refines coherence by damping unstable or high-cost spectral modes.\par
\noindent\textbf{3. Correlation Kernel}\par
\[
G(x)= e^{-\lvert x\rvert/\xi}\cos(k_F x)
\]
\noindent This closes the progression by turning matrix coherence into a physical correlation length on the lattice.\par

\clearpage
\noindent\textbf{[ID 0106] Phase-Coherence Matrix}\par
\vspace{0.35em}
\tzpsection{Dictionary}
\begingroup\small
\tzpsection{Definition}
Phase-Coherence Matrix is a Operator. Define the functional Synchronize(helix, torus) = max \{ Re( helix | phaseCoherence | torus ) + topologicalEntanglement(helix, torus) + RelativeEntropy(helix || torus) \} subject to  helix  =  torus  = 1 and H\_helix + H\_torus = E\_TZP.

% BEGIN TZP CONTEXTUAL MEANING
\tzpsection{Contextual Meaning}
\begingroup\footnotesize\setlength{\emergencystretch}{3em}\sloppy
\textbf{Formal statement:} Define the functional Synchronize(helix, torus) = max \{ Re( helix | phaseCoherence | torus ) + topologicalEntanglement(helix, torus) + RelativeEntropy(helix || torus) \} subject to  helix  =  torus  = 1 and H\_helix + H\_torus = E\_TZP.
\textbf{Contextual meaning:} This matrix encodes how phases on a helical lattice are correlated. Modes with larger Laplacian eigenvalues are exponentially suppressed, yielding finite correlation lengths.
\textbf{Derivable consequences:} Equilibrium conditions for helical toroidal coupling; criteria for energy extraction; relations between coherence length and entanglement; links to topological entropy.
\textbf{Lemma support:} Define the functional Synchronize(helix, torus) = max \{ Re( helix | phaseCoherence | torus ) + topologicalEntanglement(helix, torus) + RelativeEntropy(helix || torus) \} subject to  helix  =  torus  = 1 and H\_helix + H\_torus = E\_TZP.\par
\endgroup
% END TZP CONTEXTUAL MEANING

\tzpsection{Technical Interpretation}
Compute the functional for candidate modes and compare with observed helical--toroidal couplings; test if nature selects configurations maximizing this functional. ID 106: Phase Coherence Matrix

\tzpsection{Equation Grounding}
Define the functional Synchronize(helix, torus) = max \{ Re( helix | phaseCoherence | torus ) + topologicalEntanglement(helix, torus) + RelativeEntropy(helix || torus) \} subject to  helix  =  torus  = 1 and H\_helix + H\_torus = E\_TZP.

\tzpsection{Interpretive Role}
This matrix encodes how phases on a helical lattice are correlated. Modes with larger Laplacian eigenvalues are exponentially suppressed, yielding finite correlation lengths.
\endgroup

\tzpsection{TZP Thesaurus}
\begin{enumerate}[leftmargin=1.6em,itemsep=6pt,topsep=4pt]
\item \textbf{Pairwise Coherence Entry}: The introduction of synchronization as a measurable comparison between phase sectors.
\item \textbf{Spectral Weighting}: The refinement of coherence by damping unstable or high-cost spectral modes.
\item \textbf{Correlation Kernel}: The translation of matrix coherence into a physical correlation length on the lattice.
\end{enumerate}

\tzpsection{Encyclopedia Mathematica}
\begingroup\footnotesize\sloppy
The visual plate below is reserved for the source-specific equation and progression graphic for [ID 0106]. It is included only as a companion to the canonical equation, not as a substitute for the formal text.
\par\endgroup
\ifTZPDarkEdition
\tzpdictionaryvisualband{D:/TZPID/kdp_workflow/build/tikz_figures/ID0106_dictionary_figure_dark.pdf}
\fi

\clearpage
\begingroup\tzpsupportsetup
\noindent\textbf{\fontsize{10.8pt}{12.8pt}\selectfont [ID 0106] Constructive Logic and Proof Certificate}\par
\noindent\textbf{\fontsize{10pt}{12pt}\selectfont Phase-Coherence Matrix}\par
\vspace{0.45em}
\begingroup
\footnotesize
\setlength{\parskip}{0.22em}
\setlength{\parindent}{0pt}
\noindent\textbf{Curry--Howard--Lambek Interpretation}\par
Under the Curry--Howard--Lambek correspondence, this registry entry is read as a typed proof
object: the stated source condition supplies the proposition, the derivation supplies the
morphism, and the proof certificate records the machine handles needed to check the obligation
in the formal registry.

\vspace{0.45em}
\noindent\textbf{CHL Proof}\par
\textbf{Statement:} this matrix encodes how phases on a helical lattice are correlated

\textbf{Logic Validation:}\\
\textbf{Proposition:} ``Phase-Coherence Matrix is valid when the stated geometric or topological relation
preserves its invariant.''\\
\textbf{Proof:} The source statement names a geometric/topological carrier. The proof obligation is to
show that deformation, projection, or passage through the stated structure leaves the
relevant invariant or admissible region well defined.

\textbf{VERDICT:} \ensuremath{\checkmark}\ \textbf{TRUE} --- topological-invariance validation

\textbf{Formal Status:} \ensuremath{\checkmark}\ THEORETICAL -- source-backed CHL proof obligation

\textbf{Physical Status:} THEORETICAL -- requires model-specific empirical interpretation
\par\vspace{0.45em}
\noindent{\tiny\textbf{Isabelle/HOL.} \tzpplaincode{physics\_validity\_from\_model, external\_validity\_package, registry\_number\_matches, equation\_count\_matches}}\par
\ifTZPDarkEdition
\tzpproofvisualband{D:/TZPID/kdp_workflow/build/tikz_figures/ID0106_proof_certificate_figure_dark.pdf}
\fi
\endgroup
\endgroup
